\section{Wire Format Engine and Optimization}

\subsection{Encoding Primitives}
The serialization engine in \texttt{ctpb::wire} is designed for direct byte-stream manipulation without intermediate allocation.

\subsubsection{Varint Encoding}
Varints represent values using 7 bits of data per byte, with the most significant bit (MSB) indicating whether additional bytes follow:
\begin{lstlisting}[language=ModernC++, caption={Varint Encoding Routine}]
inline void write_varint(std::vector<std::byte>& out, std::uint64_t v) {
    while (v >= 0x80) {
        out.push_back(static_cast<std::byte>((v & 0x7F) | 0x80));
        v >>= 7;
    }
    out.push_back(static_cast<std::byte>(v & 0x7F));
}
\end{lstlisting}

\subsubsection{Zigzag Arithmetic}
Signed integers are converted into unsigned representations via single-instruction bit shifts and XORs (\texttt{write\_zigzag32}, \texttt{write\_zigzag64}), preventing small negative numbers from consuming 10 bytes on the wire.

\subsubsection{Packed Repeated Fields}
In proto3, repeated numeric scalar fields are packed by default into a single length-delimited record. \texttt{ctpb} computes the packed payload size inline and writes the contiguous sequence with a single tag.

\subsection{Deserialization and Error Recovery}
The \texttt{ctpb::wire::reader} class operates on a non-owning \texttt{std::span<const std::byte>}. All operations enforce bounds checking:
\begin{lstlisting}[language=ModernC++, caption={Bounds-Checked Varint Decoding}]
constexpr bool reader::read_varint(std::uint64_t& out) noexcept {
    std::uint64_t result = 0;
    int shift = 0;
    while (pos_ < buf_.size() && shift < 70) {
        std::byte b = buf_[pos_++];
        result |= static_cast<std::uint64_t>(
            static_cast<std::uint8_t>(b) & 0x7F) << shift;
        if ((static_cast<std::uint8_t>(b) & 0x80) == 0) {
            out = result;
            return true;
        }
        shift += 7;
    }
    return false;
}
\end{lstlisting}

When unknown fields are encountered, the \texttt{skip\_field} reader method uses the wire type to safely advance past unexpected integers, fixed-size blocks, or length-delimited payloads, maintaining forward compatibility with newer schema versions.
